\subsection{Short-curve processes and the systole}
\label{sec:pants-spectrum}

We now determine the limiting process of short curves and the systole.
The moment estimates proved here will also be used in
Subsection~\ref{sec:pants-eigenvalues} to count many disjoint short pairs of pants.
Throughout this and the next subsection,
$g\to\infty$, $n/\sqrt g\to\infty$, and $X$ has law
$\mathbb P^{\mathrm{WP}}_{g,n}$. Let $\mathcal P_2(X)$ be the set of
unoriented simple closed geodesics bounding a pair of pants with two
cusps. Define the deterministic intensity parameter
\begin{equation}
 q_g=\frac{n(n-1)}{8g\mathcal C(n/g)}.
 \label{eq:PS1}
\end{equation}
The profile endpoints in Theorem~\ref{thm:volume} imply
$q_g\asymp n^2/(g+n)\to\infty$. In particular
$q_g\sim n^2/(32\pi^2g)$ when $n/g\to0$, and
$q_g\sim x_cn/(8\pi^2)$ when $n/g\to\infty$.

We write $\operatorname{PPP}(\nu)$ for a Poisson point process with intensity measure $\nu$.

\begin{theorem}[Short-curve processes and the systole]\label{thm:pants-poisson}
\emph{(i) Curves cutting off two cusps.} The random locally finite counting measures
\begin{equation}
 \sum_{\gamma\in\mathcal P_2(X)}
       \delta_{\sqrt{q_g}\ell_\gamma}
 \ \Longrightarrow\ \operatorname{PPP}(2t\,\mathrm{d}t)
 \quad\hbox{on }[0,\infty)
 \label{eq:PS2}
\end{equation}
converge in distribution for the vague topology. Thus counts in finitely
many disjoint bounded intervals $[a_i,b_i]$ converge jointly to independent
Poisson variables with means $b_i^2-a_i^2$. If $\ell_{\min}^{(2)}$ is the
shortest length in $\mathcal P_2(X)$, with value infinity when this set is
empty, then
\[
 \mathbb P^{\mathrm{WP}}_{g,n}
 \{\sqrt{q_g}\ell_{\min}^{(2)}>t\}\longrightarrow \mathrm{e}^{-t^2}
 \qquad(t\ge0).
\]
\smallskip
\noindent\emph{(ii) All planar cusp types.} Let $\mathcal P_k(X)$ consist of simple geodesics whose
genus-zero side contains exactly $k$ cusps. If $n/g\to\lambda\in[0,\infty]$,
put $x_\lambda=X(Y_\lambda)$, where $\lambda(Y_\lambda)=\lambda$, with
$x_0=0$ and $x_\infty=x_c$. Define
\begin{equation}
 w_k(x)=\frac{2V_{0,k+1}}{k!}
             \left(\frac{x}{2\pi^2}\right)^{k-2},\qquad
 W(x)=\sum_{k\ge2}w_k(x)=\frac{2\Phi_0'(x)}{x^2}.
 \label{eq:PC1}
\end{equation}
Here $w_2(0)=W(0)=1$, $w_k(0)=0$ for $k>2$, and $W(x_c)<\infty$.
Then
\begin{equation}
 \sum_{k\ge2}\ \sum_{\gamma\in\mathcal P_k(X)}
       \delta_{(\sqrt{q_g}\ell_\gamma,k)}
 \ \Longrightarrow\
 \operatorname{PPP}\!\left(2t\,\mathrm{d}t\otimes
                    \sum_{k\ge2}w_k(x_\lambda)\delta_k\right).
 \label{eq:PC2}
\end{equation}
Convergence is for the vague topology on the product with the discrete
type space. The projection to lengths also converges, to
$\operatorname{PPP}(2W(x_\lambda)t\,\mathrm{d}t)$.

\noindent\emph{(iii) Primitive geodesics and the systole.} This is also the limit of
the process of all unoriented primitive closed geodesics, at the same
rescaling. Writing $\operatorname{sys}(X)$ for the shortest closed
geodesic length, we have
\[
 \mathbb P^{\mathrm{WP}}_{g,n}
 \{\sqrt{q_g}\operatorname{sys}(X)>t\}
 \longrightarrow \mathrm{e}^{-W(x_\lambda)t^2}\qquad(t\ge0).
\]
With probability tending to one every systolic curve is planar
separating. Select one by any measurable tie rule, and let $K_g$ be
its planar-side cusp count, with value zero if it is not planar. Then
$\bigl(\sqrt{q_g}\operatorname{sys}(X),K_g\bigr)$ converges jointly
to a pair $(Z,K)$ with
\begin{equation}
 \begin{aligned}
 \mathbb P\{Z\in \mathrm{d}t,K=k\}
 &=2t\,w_k(x_\lambda)\mathrm{e}^{-W(x_\lambda)t^2}\,\mathrm{d}t,\\
 \mathbb P\{K=k\}&=\frac{w_k(x_\lambda)}{W(x_\lambda)}
 \qquad(t>0,\ k\ge2).
 \end{aligned}
 \label{eq:FS1}
\end{equation}
\end{theorem}

The method of factorial moments for curves separating off a prescribed
number of cusps has a direct precedent in Hide--Thomas~\cite[Theorem~1.7]{HTSphere}
for punctured spheres, which also treats summed types and nested
configurations~\cite[Remark~1.9 and Proposition~3.6]{HTSphere}. Here the uniform volume ratios
give the entire family of type intensities for diverging genus at every
cusp density above the square-root scale. These planar curves
become contractible after the cusps are filled; they are excluded from
the tight spectrum of~\cite[Definition~2.1]{BL}. The all-primitive
process of Mirzakhani--Petri~\cite[Theorem~4.1]{MP} concerns closed
surfaces at fixed physical lengths. Here the full process is observed
on a shrinking scale; a tight-spectrum limit at its longer natural
scale is a separate question.

\emph{Factorial moments.} For sufficiently small $L$, the collar lemma forces all curves counted by
$N_L=\#\{\gamma\in\mathcal P_2(X):\ell_\gamma\le L\}$ to be disjoint.
Their pants interiors are also disjoint. Otherwise one pants boundary
would be an interior essential geodesic of the other pants, hence
boundary-parallel, contradicting uniqueness of geodesic representatives.
Thus their pairs of cusp labels are disjoint, and their complement has
type $(g,n-r)$ for $r$ pants. We have $g\ge2$ eventually, so the
exceptional double-pants decomposition of a four-punctured sphere does
not arise.

Mirzakhani's integration formula~\cite{MirInv,Mir}, with the same conventions as
\cite[Lemma~29]{SW}, gives, for any fixed $r$ and short length windows
$I_1,\ldots,I_r$, the ordered distinct-curve integral
\begin{equation}
 \mathbb E\!\sum_{\substack{\gamma_1,\ldots,\gamma_r\in\mathcal P_2(X)\\
                            \text{distinct}}}
       \prod_{j=1}^r\mathbf1_{\ell_{\gamma_j}\in I_j}
 =\frac{(n)_{2r}}{2^r V_{g,n}}
    \int_{I_1\times\cdots\times I_r}
       V_{g,n-r}(x_1,\ldots,x_r,0^{n-2r})
       \prod_{j=1}^r x_j\,\mathrm{d}x_j.
 \label{eq:PS3}
\end{equation}
Here $(n)_{2r}=n(n-1)\cdots(n-2r+1)$. The factor $2^{-r}$ removes
the order of the two labels within each pair. The tuple of curves is
ordered, so there is no $r!$ divisor. Each pants volume is $V_{0,3}=1$.
Mirzakhani's boundary estimate~\cite[Lemma~3.2]{Mir}, in the form
of~\cite[Lemma~14]{SW},
\begin{equation}
 1\le\frac{V_{g,m}(\mathbf x)}{V_{g,m}}
       \le\exp\!\left(\frac12\sum_jx_j\right)
 \label{eq:PS4}
\end{equation}
is uniform in the topological type. For fixed $r$, it replaces the
boundary volume in \eqref{eq:PS3} by $V_{g,n-r}(1+O_r(L))$ when
all $I_j\subset[0,L]$.

Theorem~\ref{thm:volume} also implies, for every fixed $r$,
\begin{equation}
 \frac{V_{g,n-r}}{V_{g,n}}
  =(2g\mathcal C(n/g))^{-r}(1+o(1)).
 \label{eq:PS5}
\end{equation}
Indeed this is a product of $r$ adjacent cusp ratios. Replacing $(n-j)/g$
by $n/g$ changes their profile by $1+O_r(n^{-1})$, because
$\lambda\partial_\lambda\log\mathcal C(\lambda)$ is bounded: this follows
from analyticity in the interior and the two endpoint expansions.

\begin{proof}[Proof of Theorem~\ref{thm:pants-poisson}]
We prove the planar process by bounding the tail over cusp types and
computing the mixed factorial moments for any fixed collection of types.
We then exclude the other topological types on the same length scale;
this identifies the full primitive process and the systole. Write $P=2\pi^2$, $h=2g-1$, $m=n-1$, and $F_h=\Phi_g'$.
Choose $x\in(0,x_c)$ by $h\Psi(x)=m$ and put $q_g^*=nx/(4P)$.
The effective density is $2(n-1)/(2g-1)$. The boundedness of
$\lambda\partial_\lambda\log X(Y_\lambda)$, including its endpoint limits, implies
$q_g^*/q_g\to1$ and $q_g^*\to\infty$.

\emph{Step 1: a uniform tail bound over cusp types.}
We first show that curves cutting off a large number of cusps contribute
negligibly, uniformly along the sequence. The input is the
coefficient-ratio bound \eqref{eq:PC3}, which compares volumes whose cusp
counts differ by $j$.
By the positive convolution \eqref{eq:TL33} and
Lemma~\ref{lem:point-mass}, every point mass is at most $C/\sigma$,
where $\sigma^2=h\mathscr D\Psi(x)$. The local limit
\eqref{eq:TL47} gives a lower bound $c/\sigma$ at $m$.
Consequently, for all sufficiently large $h,m$,
\begin{equation}
 \frac{[z^{m-j}]F_h(z)}{[z^m]F_h(z)}
 =x^j\frac{\mathbb P_x(N=m-j)}{\mathbb P_x(N=m)}
 \le Cx^j\qquad(0\le j\le m).
 \label{eq:PC3}
\end{equation}
Only the denominator uses the local limit; the numerator may be at any
index, including zero.

Let $N_{k,L}$ count curves of type $k$ with length at most $L$, and put
$a_k=V_{0,k+1}/(k!P^{k-2})$. Mirzakhani integration gives
\begin{equation}
 \begin{aligned}
 \mathbb E N_{k,L}=\frac{\binom nk}{V_{g,n}}
 \int_0^L\ell\,V_{0,k+1}(\ell,0^k)
                  V_{g,n-k+1}(\ell,0^{n-k})\,\mathrm{d}\ell,\\
 \binom nk\frac{V_{0,k+1}V_{g,n-k+1}}{V_{g,n}}
 =\frac nP a_k\frac{[z^{n-k}]F_h}{[z^{n-1}]F_h}.
 \end{aligned}
 \label{eq:PC4}
\end{equation}
The two sides have different genera, so there is no side-exchange
factor. The coefficient identity follows from \eqref{eq:TL50}; it
holds up to $k=n$. Bound \eqref{eq:PS4} changes the integrand's two
volumes by a factor between one and $\mathrm{e}^L$, independently of $k$.
Using \eqref{eq:PC3} with $j=k-1$ therefore gives, for each fixed $T$,
\begin{equation}
 \mathbb E\sum_{k\ge K}N_{k,T/\sqrt{q_g^*}}
 \le CT^2\sum_{k\ge K}a_kx_c^{k-2}\longrightarrow0
 \quad(K\to\infty),
 \label{eq:PC5}
\end{equation}
uniformly along the sequence. To justify summability directly, the
genus-zero string solution gives $\Phi_0'''(x)=1/a(Y(x))=Y'(x)$.
Since $V_{0,3}=1$, its zero initial conditions give
$\Phi_0''=Y$ and $\Phi_0'(x)=\int_0^xY(u)\,\mathrm{d}u$. This is finite at
$x_c$; positivity gives convergence of the coefficient series there.
Its initial term $x^2/2$ also gives $W(0)=1$.

\emph{Step 2: fixed types and nested configurations.}
Next we compute the mixed factorial moments for finitely many fixed cusp
types. The main contribution comes from configurations whose planar
sides are pairwise disjoint, which we call \emph{stars}; nested
configurations are shown to be negligible.
Fix $r$ ordered distinct curves of prescribed types $k_1,\ldots,k_r$,
and write $K=\sum k_i$. The collar lemma makes their boundaries
disjoint. For a star, the ordered
cusp-label multiplicity is $(n)_K/\prod_i k_i!$, and the complement
has type $(g,n-K+r)$. There is no $r!$ divisor. The boundary error is
$1+O_r(L)$, as each cut occurs twice in \eqref{eq:PS4}.
For fixed $d$, \eqref{eq:PS5} and $q_g^*/q_g\to1$ give
$V_{g,n-d}/V_{g,n}=(x/(Pn))^d(1+o(1))$.
Thus for windows $[a_i,b_i]/\sqrt{q_g^*}$ the star contribution tends to
\begin{equation}
 \prod_{i=1}^r
       w_{k_i}(x_\lambda)(b_i^2-a_i^2).
 \label{eq:PC6}
\end{equation}
The positive expression has a relative error before taking the limit,
even when some limiting weights vanish.

Filling cusps turns each planar side into a disk. Disjoint boundaries
therefore have nested or disjoint disks; their cut graph is a tree with
one genus-$g$ root and $r$ planar vertices. Let $d$ be the root degree,
$a_v$ the number of cusps at a planar vertex and $d_v$ its degree;
put $a=\sum_v a_v$. Prescribed types bound $a$ by $K$.
There are finitely many such trees with ordered edges and cusp
allocations, depending only on the fixed types. Classification of
surfaces gives the corresponding mapping-class orbits with bounded
symmetry factors; assigning cusp labels costs at most $Cn^a$.
The root volume ratio is $O((x/(Pn))^{a-d})$, the planar volumes
are fixed, and the $r$ length integrals cost $O((nx)^{-r})$.
Each cut type contributes at most
\begin{equation}
 Cn^{d-r}x^{a-d-r}=O(n^{-1})\quad\hbox{if it is not a star}.
 \label{eq:PC7}
\end{equation}
Indeed stability gives $a_v+d_v\ge3$ and
$\sum_vd_v=2r-d$, hence $a\ge r+d$. The power of $x$ is
nonnegative, so no loss occurs as $x\to0$. A nonstar has $d\le r-1$;
if $d=r$, every planar vertex is a leaf and the configuration is a star.
This accounts for every fixed-type mixed factorial moment.

The limits \eqref{eq:PC6} are the mixed factorial moments of independent
Poisson counts. Higher fixed moments give uniform integrability, and
the product Poisson law is moment determinate. First moments give
tightness on compact sets in the marked space. This proves
\eqref{eq:PC2}. Its unmarked projection requires more than vague
convergence: truncate the types, pass to the finite-type limit, then
let the cutoff grow using \eqref{eq:PC5} and the convergent limiting
series. This proves the stated unmarked process. Finally replace
$q_g^*$ by $q_g$. The type-two restriction gives \eqref{eq:PS2}
without assuming a limit for $n/g$, by subsequences and compactness of
$[0,\infty]$. Its zero-count probabilities give the pants minimum law.

\emph{Step 3: the remaining curves and the systole.}
It remains to show that, on the scale $q_g^{-1/2}$, there are no short
curves other than planar separating ones: nonseparating curves, and
separating curves with positive genus on both sides, are rare.
Nonseparating simple curves of length at most $L\le L_0$ have expected
count at most $C\mathrm{e}^LL^2V_{g-1,n+2}/V_{g,n}\le CL^2$, by integration,
\eqref{eq:PS4}, and the uniform lower bound for $B_{g,n}$.
To sum every positive-genus separating cut, put
$S_g(z)=\sum_{i=1}^{g-1}\Phi_i'(z)\Phi_{g-i}'(z)$. Exact positivity
\eqref{eq:TL33} includes genus one, so the child powers add to $h-1$:
\begin{equation}
 S_g(z)=V(z)^{h-1}Q_g(z),\qquad Q_g\ge0
 \ \hbox{coefficientwise},\qquad
 \frac{S_g(x)}{F_h(x)}\le\frac{Ca(Y)r(Y)^2}{g^2}.
 \label{eq:FS2}
\end{equation}
The last bound is \eqref{eq:TL26}: undoing \eqref{eq:TL9} leaves
exactly $ar^2$, as the child complexities sum to $2g-2$ and their
deficits to $3g-4$. In particular bounded child genera are retained.
Lemma~\ref{lem:point-mass}, applied to the positive factor $V^{h-1}$,
gives $\sup_j\mathbb P_x^{S}(N=j)\le C/\sigma$. Against the parent lower
bound $\mathbb P_x^{F}(N=n-1)\ge c/\sigma$, the standard deviations
cancel. The coefficient identity \eqref{eq:TL50} therefore yields
\begin{equation}
 \begin{aligned}
 \mathcal A_{g,n}
 &:=\sum_{i=1}^{g-1}\sum_{k=0}^n\binom nk
       \frac{V_{i,k+1}V_{g-i,n-k+1}}{V_{g,n}}
   =\frac1P\frac{[z^n]S_g(z)}{[z^n]\Phi_g(z)}\\
 &=\frac{n}{Px}\frac{S_g(x)}{F_h(x)}
       \frac{\mathbb P_x^S(N=n)}{\mathbb P_x^F(N=n-1)}
   \le\frac{Cnar^2}{xg^2}\le\frac Cg.
 \end{aligned}
 \label{eq:FS3}
\end{equation}
Indeed $(n-1)/x=h/(Aar)$, while $A$ is bounded below and $r$ above.
Thus no small-density factor is lost. Each separating orbit with two
positive genera occurs in this ordered sum. Its integration symmetry
and one-handle factors are at most one; applying \eqref{eq:PS4} to
both components bounds its total expected count by
$C\mathrm{e}^LL^2\mathcal A_{g,n}\le CL^2/g$. This includes all cusp partitions.

For fixed $T$, take $L=T/\sqrt{q_g}$. Markov's inequality shows that
with probability $1-O(T^2/q_g)$ every simple geodesic in this window
is planar. Primitive geodesics shorter than $2\operatorname{arcsinh}1$
are simple, by the standard collar theorem used in
\cite[Section~1.1]{HTSphere}. Since $q_g\to\infty$, the full primitive
process and planar process agree on each bounded rescaled window
with probability tending to one. This transfers the process limit
and its zero-count law to the systole.

Finally restrict the first point to $[\eta,T]$ and finitely many cusp
types. The first-point map is continuous when its first length is
unique and avoids the endpoints, as happens almost surely in the
limiting Poisson process. Let the type cutoff grow using
\eqref{eq:PC5}, then $\eta\downarrow0$ and $T\to\infty$.
The omitted length probabilities tend to zero by the zero-count law;
$W\ge1$. Nonplanar competitors are excluded on each fixed window.
This also makes any tie rule immaterial and proves planarity of every
systolic curve with probability tending to one. The first-arrival
density of the limiting marked process is \eqref{eq:FS1}.
\end{proof}

\subsection{Many small eigenvalues}\label{sec:pants-eigenvalues}

\begin{proposition}[Many short pants]\label{prop:pants-concentration}
If $L=L_g\to0$ and $q_gL_g^2\to\infty$, then
\begin{equation}
 \frac{N_{L_g}}{q_gL_g^2}\longrightarrow1
 \quad\hbox{in probability}.
 \label{eq:PS6}
\end{equation}
Consequently, for each fixed $\varepsilon\in(0,1)$, with probability
tending to one there are at least
$\lfloor(1-\varepsilon)q_gL_g^2\rfloor-1$ positive discrete eigenvalues
in $(0,6L_g]$, counted with multiplicity.
\end{proposition}
\begin{proof}
The first two cases of \eqref{eq:PS3} give
\[
 m_L:=\mathbb E N_L=\frac{(n)_2L^2}{4}
       \frac{V_{g,n-1}}{V_{g,n}}(1+O(L))\sim q_gL^2,
\]
and
\[
 \frac{\mathbb E(N_L)_2}{m_L^2}
 =\frac{(n-2)(n-3)}{n(n-1)}D_{g,n-1}^{-1}(1+O(L))\longrightarrow1.
\]
Since $m_L\to\infty$, the relative variance tends to zero, proving
\eqref{eq:PS6}. This second-moment argument does not use the point-process
limit at a window growing with $g$.

Hide--Thomas~\cite[Lemma~5.3]{HTSphere} prove that $N$ such pants
boundaries of length at most $L$ give $N-1$ positive eigenvalues below
$6L$, provided $6L<1/4$. For the reader's convenience we recall the
variational argument.
Inside each pants choose a function equal to inward distance $\rho$ on
the half-collar $0\le\rho\le1$, equal to one beyond it, and zero outside
the pants. For all sufficiently small boundary lengths $\ell$, the
collar lemma supplies this embedded half-collar. The function lies in
$H^1(X)$, and
\[
 \int|\nabla f|^2=\ell\sinh1,\qquad
 \int|f|^2\ge2\pi-\ell\sinh1.
\]
Different pants have disjoint supports up to their boundaries, so the
span of these functions has dimension $N$ and Rayleigh quotient at most
$6L$ for small $L$. Its intersection with the orthogonal complement of
constants has dimension at least $N-1$. Since the essential spectrum
starts at $1/4$, the form version of min--max gives the claimed positive
discrete eigenvalues. Apply \eqref{eq:PS6} and let $g$ grow.
\end{proof}

\begin{corollary}[A growing number of small eigenvalues]\label{cor:many-spectrum}
Suppose $g\to\infty$ and $n/\sqrt g\to\infty$. If $j=j(g)\to\infty$
and $j=o(q_g)$, then
\[
 \mathbb P^{\mathrm{WP}}_{g,n}
 \{\lambda_j(X)\text{ exists and }\lambda_j(X)\le12\sqrt{j/q_g}\}
 \longrightarrow1,
\]
where positive discrete eigenvalues are ordered with multiplicity.
\end{corollary}

\begin{proof}[Proof of Corollary~\ref{cor:many-spectrum}]
For $j\to\infty$ with $j=o(q_g)$, choose $L_g=2\sqrt{j/q_g}$.
Then $L_g\to0$, $q_gL_g^2=4j$, and \eqref{eq:PS6} implies
$N_{L_g}\ge j+1$ with probability tending to one. The preceding
variational estimate gives $\lambda_j\le6L_g=12\sqrt{j/q_g}$.
\end{proof}

The spectral mechanism is established in the fixed-genus work of
Hide--Thomas~\cite[Theorem~1.10 and Remark~5.5]{HTSphere}; their later
bordered-volume asymptotics give linear counts below each fixed positive
threshold at fixed genus~\cite[Theorem~1.1]{HTLargeN}.
Their deterministic result~\cite[Theorem~1.3]{HTMany}, in the July 2026
revision, gives a linear number below $1/4$ whenever $g<an$ for fixed
$a>0$. Our threshold tends to zero and the genus grows. In the range
$\sqrt g\ll n\ll g$, the corresponding second-moment volume inputs
already follow from Shen--Wu~\cite{SW}. The extension supplied here is
the density-dependent intensity and removal of the upper restriction on
$n/g$. No sharp eigenvalue asymptotic, matching lower bound, or
Poisson eigenvalue law is asserted.
